% Source: https://artofproblemsolving.com/community/c3h1877218_probability_that_two_children_like_peas_problem
% Stable identity: aops-topic-1877218-q01
(a) By independence the probability is $0.04\cdot0.04=\boxed{0.0016}$.

(b) There are $\dbinom{100}{2}=4950$ equally likely pairs, of which $\dbinom42=6$ consist of two pea-lovers: $\dfrac{6}{4950}=\boxed{\dfrac{1}{825}}\approx0.00121$.

(c) With $L=0.04N$ pea-lovers,
\[
p_N=\dfrac{\dbinom L2}{\dbinom N2}=\dfrac{L(L-1)}{N(N-1)}=0.04\cdot\dfrac{0.04N-1}{N-1},
\]
which depends on $N$ (for $N=200$ it is $\dfrac{8\cdot7}{200\cdot199}=\dfrac{7}{4975}\approx0.00141$). As $N\to\infty$, $\dfrac{0.04N-1}{N-1}\to0.04$, so $p_N\to0.04^2=0.0016$.

The two answers use different models. The original question gives a probability for \emph{each} child, with no fixed class, which is model (a), so the intended answer is $0.0016$. Model (b) samples without replacement from a fixed finite group in which exactly $4\%$ like peas; choosing one pea-lover slightly lowers the chance that the second child is one, which is why its answer depends on the class size and approaches $0.0016$ for large classes.
